% GENERATED FILE. Edit canonical lessons, then run npm run generate:books.
% canonical-source-hash: fnv1a64:e501e2358a0f8868
% canonical-lessons: ML-C276-eppolum, ML-C276-pativayi, ML-C276-palappolum, ML-C276-mikkappolum, ML-C276-apurvamayi

\chapter{Always, Regularly, Often, Usually, Rarely}
\label{ch:ml-always-regularly-often-usually-rarely}
\hlchaptermodality{276}

\begin{chapteropening}
\textbf{By the end of this chapter:} \emph{I can say always, regularly, often, usually and rarely.}

Complete the last lesson of chapter 276: I can say always, regularly, often, usually and rarely.
\end{chapteropening}

\section[\texorpdfstring{\ml{എപ്പോഴും} (eppōḻuṁ)}{എപ്പോഴും (eppōḻuṁ)}]{\ml{എപ്പോഴും} (eppōḻuṁ) — always}

\label{lesson:ML-C276-eppolum}



\begin{quote}
Before the new one: say the Malayalam for jaundice, then the Malayalam for blood pressure.
\end{quote}

\subsection*{You'll want to know: \ml{എപ്പോഴും}}
\textbf{\ml{എപ്പോഴും}} — \emph{eppōḻuṁ} — \textquotedblleft{}always\textquotedblright{}.

\begin{cousinweb}
\textbf{\ml{എപ്പോഴും}} (\emph{eppōḻuṁ}) is native Dravidian: \emph{eppōḷ}, \textquotedblleft{}when\textquotedblright{}, which DEDR 5151 lists, built on the word for \textquotedblleft{}time\textquotedblright{} (DEDR 4559), with \emph{-uṁ}: \textquotedblleft{}at any time\textquotedblright{}.
\end{cousinweb}

\begin{grammarlens}[title={What you've built}]
One more word for this chapter.
\end{grammarlens}

\subsection*{Your turn}
\begin{itemize}
\raggedright
  \item \emph{Say it:} \emph{eppōḻuṁ}
  \item \emph{Say it:} \emph{eppōḻuṁ}, once more
  \item \emph{Say it:} say \emph{eppōḻuṁ}
  \item \emph{Recall:} say the Malayalam for jaundice, then the Malayalam for blood pressure, then say \emph{eppōḻuṁ} again
\end{itemize}

\begin{tcolorbox}[breakable,colback=teal!4,colframe=teal!35!black,title={Before you move on}]
What does \ml{എപ്പോഴും} mean? (\textquotedblleft{}Always\textquotedblright{}.) Say it once more.
\end{tcolorbox}
\section[\texorpdfstring{\ml{പതിവായി} (pativāyi)}{പതിവായി (pativāyi)}]{\ml{പതിവായി} (pativāyi) — regularly}

\label{lesson:ML-C276-pativayi}



\begin{quote}
Before the new one: say the Malayalam for blood pressure, then the Malayalam for always.
\end{quote}

\subsection*{You'll want to know: \ml{പതിവായി}}
\textbf{\ml{പതിവായി}} — \emph{pativāyi} — \textquotedblleft{}regularly\textquotedblright{}.

\begin{cousinweb}
\textbf{\ml{പതിവായി}} (\emph{pativāyi}) is native Dravidian: \emph{pativu}, \textquotedblleft{}settlement, custom\textquotedblright{} (DEDR 3911, with Tamil \emph{pati}), and \emph{āyi}, \textquotedblleft{}as, having become\textquotedblright{} (DEDR 333): \textquotedblleft{}as a custom\textquotedblright{}. DEDR 333 compares Sanskrit \emph{ām}.
\end{cousinweb}

\begin{grammarlens}[title={What you've built}]
One more word for this chapter.
\end{grammarlens}

\subsection*{Your turn}
\begin{itemize}
\raggedright
  \item \emph{Say it:} \emph{pativāyi}
  \item \emph{Say it:} \emph{pativāyi}, once more
  \item \emph{Say it:} say \emph{pativāyi}
  \item \emph{Recall:} say the Malayalam for blood pressure, then the Malayalam for always, then say \emph{pativāyi} again
\end{itemize}

\begin{tcolorbox}[breakable,colback=teal!4,colframe=teal!35!black,title={Before you move on}]
What does \ml{പതിവായി} mean? (\textquotedblleft{}Regularly\textquotedblright{}.) Say it once more.
\end{tcolorbox}
\section[\texorpdfstring{\ml{പലപ്പോഴും} (palappōḻuṁ)}{പലപ്പോഴും (palappōḻuṁ)}]{\ml{പലപ്പോഴും} (palappōḻuṁ) — often}

\label{lesson:ML-C276-palappolum}



\begin{quote}
Before the new one: say the Malayalam for always, then the Malayalam for regularly.
\end{quote}

\subsection*{You'll want to know: \ml{പലപ്പോഴും}}
\textbf{\ml{പലപ്പോഴും}} — \emph{palappōḻuṁ} — \textquotedblleft{}often\textquotedblright{}.

\begin{cousinweb}
\textbf{\ml{പലപ്പോഴും}} (\emph{palappōḻuṁ}) is native Dravidian, \textquotedblleft{}many times\textquotedblright{}: \emph{pala}, \textquotedblleft{}many\textquotedblright{} (DEDR 3987, with Tamil \emph{pala}), the word for \textquotedblleft{}time\textquotedblright{} (DEDR 4559), and \emph{-uṁ}.
\end{cousinweb}

\begin{grammarlens}[title={What you've built}]
One more word for this chapter.
\end{grammarlens}

\subsection*{Your turn}
\begin{itemize}
\raggedright
  \item \emph{Say it:} \emph{palappōḻuṁ}
  \item \emph{Say it:} \emph{palappōḻuṁ}, once more
  \item \emph{Say it:} say \emph{palappōḻuṁ}
  \item \emph{Recall:} say the Malayalam for always, then the Malayalam for regularly, then say \emph{palappōḻuṁ} again
\end{itemize}

\begin{tcolorbox}[breakable,colback=teal!4,colframe=teal!35!black,title={Before you move on}]
What does \ml{പലപ്പോഴും} mean? (\textquotedblleft{}Often\textquotedblright{}.) Say it once more.
\end{tcolorbox}
\section[\texorpdfstring{\ml{മിക്കപ്പോഴും} (mikkappōḻuṁ)}{മിക്കപ്പോഴും (mikkappōḻuṁ)}]{\ml{മിക്കപ്പോഴും} (mikkappōḻuṁ) — usually, most of the time}

\label{lesson:ML-C276-mikkappolum}



\begin{quote}
Before the new one: say the Malayalam for regularly, then the Malayalam for often.
\end{quote}

\subsection*{You'll want to know: \ml{മിക്കപ്പോഴും}}
\textbf{\ml{മിക്കപ്പോഴും}} — \emph{mikkappōḻuṁ} — \textquotedblleft{}usually, most of the time\textquotedblright{}.

\begin{cousinweb}
\textbf{\ml{മിക്കപ്പോഴും}} (\emph{mikkappōḻuṁ}) is native Dravidian, \textquotedblleft{}most times\textquotedblright{}: \emph{mikka}, \textquotedblleft{}the greater part\textquotedblright{} (DEDR 4838, with Tamil \emph{miku}), the word for \textquotedblleft{}time\textquotedblright{} (DEDR 4559), and \emph{-uṁ}.
\end{cousinweb}

\begin{grammarlens}[title={What you've built}]
One more word for this chapter.
\end{grammarlens}

\subsection*{Your turn}
\begin{itemize}
\raggedright
  \item \emph{Say it:} \emph{mikkappōḻuṁ}
  \item \emph{Say it:} \emph{mikkappōḻuṁ}, once more
  \item \emph{Say it:} say \emph{mikkappōḻuṁ}
  \item \emph{Recall:} say the Malayalam for regularly, then the Malayalam for often, then say \emph{mikkappōḻuṁ} again
\end{itemize}

\begin{tcolorbox}[breakable,colback=teal!4,colframe=teal!35!black,title={Before you move on}]
What does \ml{മിക്കപ്പോഴും} mean? (\textquotedblleft{}Usually, most of the time\textquotedblright{}.) Say it once more.
\end{tcolorbox}
\section[\texorpdfstring{\ml{അപൂർവമായി} (apūrvamāyi)}{അപൂർവമായി (apūrvamāyi)}]{\ml{അപൂർവമായി} (apūrvamāyi) — rarely}

\label{lesson:ML-C276-apurvamayi}



\begin{quote}
Before the new one: say the Malayalam for often, then the Malayalam for usually.
\end{quote}

\subsection*{You'll want to know: \ml{അപൂർവമായി}}
\textbf{\ml{അപൂർവമായി}} — \emph{apūrvamāyi} — \textquotedblleft{}rarely\textquotedblright{}.

\begin{cousinweb}
\textbf{\ml{അപൂർവമായി}} (\emph{apūrvamāyi}) joins Sanskrit \emph{apūrva}, \textquotedblleft{}unprecedented, rare\textquotedblright{} (CDIAL 494), to native \emph{āyi}, \textquotedblleft{}as\textquotedblright{} (DEDR 333). DEDR 333 compares Sanskrit \emph{ām}.
\end{cousinweb}

\begin{grammarlens}[title={What you've built}]
One more word for this chapter.
\end{grammarlens}

\subsection*{Your turn}
\begin{itemize}
\raggedright
  \item \emph{Say it:} \emph{apūrvamāyi}
  \item \emph{Say it:} \emph{apūrvamāyi}, once more
  \item \emph{Say it:} say \emph{apūrvamāyi}
  \item \emph{Recall:} say the Malayalam for often, then the Malayalam for usually, then say \emph{apūrvamāyi} again
\end{itemize}

\begin{tcolorbox}[breakable,colback=teal!4,colframe=teal!35!black,title={Before you move on}]
What does \ml{അപൂർവമായി} mean? (\textquotedblleft{}Rarely\textquotedblright{}.) Say it once more.
\end{tcolorbox}
